\emph{Setting.} The fields are those of Definition \claimref{obj.field-theory.single-algebra-field}, quantised as in Proposition \claimref{field-theory.second-quantisation}; the interaction is $H_{\mathrm{int}} = e\sum J^{\mu}A_{\mu}$ of Proposition \claimref{field-theory.coupled-system-tree-level}, and photon exchange is used in the covariant form, the map $a \mapsto -a/k^{2}$ on vectors ($-g_{\mu\nu}/k^{2}$ in components), on every line. Proposition \claimref{field-theory.elementary-photon-ga}(4) shows that this equals Coulomb-gauge exchange only at tree level on conserved external currents; on internal off-shell lines of a loop the covariant propagator is an import: it needs the indefinite-form quantisation of the gauge field, Premise \claimref{synthesis.krein-form-premise-indefinite-inner-product-for-gauge-fields} (a self-adjoint BRST charge needs it, by the gauge-fixing result \claimref{synthesis.self-adjoint-brst-charge-needs-a-krein-form-and-the-tower-supplies-the-ghosts}), and the gauge independence of on-shell quantities (Fact \claimref{synthesis.imported-one-loop-integrals-and-pauli-villars-regulators}(iv)). The rest of the one-loop machinery is standard and imported (Fact \claimref{synthesis.imported-one-loop-integrals-and-pauli-villars-regulators}): Feynman parametrisation $1/(D_{1}\cdots D_{n}) = (n-1)!\int_{\mathrm{simplex}}(\sum x_{j}D_{j})^{-n}$; the Wick-rotated integrals $\int\frac{d^{4}\ell}{(2\pi)^{4}}(\ell^{2}-\Delta)^{-3} = \frac{-I_{4}}{32\pi^{2}\Delta}$ and $\int\frac{d^{4}\ell}{(2\pi)^{4}}(\ell^{2}-\Delta)^{-2} = \frac{I_{4}}{16\pi^{2}}(\ln(L/\Delta) - 1 + O(\Delta/L))$; Pauli and Villars' regulator fields, which enter with weights $c_{i}$ of both signs (fields of negative weight); on-shell external states (Peskin and Schroeder, chs. 6 and 7). The units of the Feynman rules (one per propagator, $-e$ times the unit per vertex) and of the Wick rotation are the field theory's unit $I_{4}$, numbers, including the values of the Wick-rotated integrals, being read in $\mathrm{span}\{1, I_{4}\}$; on a Clifford product they act as the right unit, collected outside the product, and clause (5) shows that this placement is forced on massive lines. The regulator scheme: Pauli-Villars fermions on the fermion loop for (2), a Pauli-Villars photon on the photon line for (3) and (4) (clause (3) is finite and needs none). Every Clifford computation below is exact and is the content of the check.

\emph{(1) Traces.} The Dirac matrices over $\mathbb{Q}(i)$ satisfy $\gamma_{\mu}\gamma_{\nu} + \gamma_{\nu}\gamma_{\mu} = 2g_{\mu\nu}$, so the map from basis blades to matrix products is the representation of $\Cl(1,3)$. Every blade except the scalar is traceless, so $\mathrm{Tr}\,X = 4\langle X\rangle_{0}$. With the chirality $\gamma_{5}$ acting as $x \mapsto I_{4}xI_{4}$ ($I_{4}$ inside the product on the left, the unit $I_{4}$ on the right; $I_{5}\gamma_{0}\gamma_{1}\gamma_{2}\gamma_{3}$ in the matrix reading, whose complex scalar is read as the central $I_{5} = I_{4}e$ of $\Cl(2,3)$, \claimref{std-algebra.complexification-is-the-odd-extension}; through $\tfrac12(1 + e)$ it acts as the right unit $xI_{4}$, related to it there and not identified), $\gamma_{5}^{2} = 1$ and $\mathrm{Tr}(\gamma_{5}X) = 4\langle I_{4}X\rangle_{0}\,I_{4} = -4x_{4}I_{4}$, the unit in front, so $\mathrm{Tr}(\tfrac12(1 \pm \gamma_{5})X) = 2(\langle X\rangle_{0} \mp x_{4}I_{4})$, which is $2\proj(X)$ with the pseudoscalar of $X$ sent to $\mp$ the unit. For four vectors, $x_{4}$ is the determinant of their components and $\langle abcd\rangle_{0} = (a\cdot b)(c\cdot d) - (a\cdot c)(b\cdot d) + (a\cdot d)(b\cdot c)$ (checked).

\emph{(2) Vacuum polarisation.} The loop is $I_{4}\Pi^{\mu\nu}(q) = -(-eI_{4})^{2}\int\frac{d^{4}k}{(2\pi)^{4}}\mathrm{Tr}[\gamma^{\mu}(k+m)\gamma^{\nu}(k'+m)\,I_{4}I_{4}]/(D D')$, the two propagator units standing on the right of the product, and the four units giving $-e^{2}$ in front. One Feynman parameter and the shift $k = \ell - xq$, $k' = \ell + (1-x)q$ give the denominator $(\ell^{2} - \Delta)^{2}$, $\Delta = m^{2} - x(1-x)Q(q)$. The numerator $4\langle\gamma^{\mu}(k+m)\gamma^{\nu}(k'+m)\rangle_{0}$ is bilinear in $(k, k')$ plus $4m^{2}g^{\mu\nu}$. Terms linear in $\ell$ integrate to zero, and $\ell^{\alpha}\ell^{\beta} \to \tfrac14 Q(\ell)g^{\alpha\beta}$; in the bilinear form this replaces $B(\ell, \ell)$ by $\tfrac14 Q(\ell)\sum_{\alpha}B(\gamma_{\alpha}, \gamma^{\alpha})$. The check computes the result as an exact tensor at five values of $x$ (the numerator is a polynomial of degree two in $x$, so three suffice) and finds
\[
-2Q(\ell)g^{\mu\nu} + 8x(1-x)\bigl(Q(q)g^{\mu\nu} - q^{\mu}q^{\nu}\bigr) + 4\Delta g^{\mu\nu}.
\]
Read as a map on vectors, $a \mapsto \sum a_{\mu}N^{\mu\nu}\gamma_{\nu}$, the middle term is $8x(1-x)(Q(q)a - (a\cdot q)q)$, which annihilates $q$: it is transverse. The $g$ part is $-2(Q(\ell) - 2\Delta)g$; after the Wick rotation ($Q(\ell) = -u$, $d^{4}\ell = I_{4}\pi^{2}u\,du$ on angular integration) it is $2I_{4}\pi^{2}g\int_{0}^{L}u(u + 2\Delta)(u + \Delta)^{-2}\,du$, and $F(L) = L - \Delta + \Delta^{2}/(L+\Delta)$ is the exact antiderivative with $F(0) = 0$ (checked as rational functions). Summed over the regulators, $\sum c_{i}F_{i}(L) = L\sum c_{i} - \sum c_{i}\Delta_{i} + \sum c_{i}\Delta_{i}^{2}/(L + \Delta_{i})$; the first two sums vanish since $\sum c_{i} = 0$ and $\sum c_{i}\Delta_{i} = \sum c_{i}M_{i}^{2} - x(1-x)Q(q)\sum c_{i} = 0$, and the last is a proper rational function of $L$, so it tends to zero (checked; with a single regulator a constant survives, the control). Hence $\Pi(a) = (Q(q)a - (a\cdot q)q)\Pi(q^{2})$ (in indices $(Q(q)g^{\mu\nu} - q^{\mu}q^{\nu})\Pi(q^{2})$) with
\[
\Pi(q^{2}) = -e^{2}\cdot\frac{1}{16\pi^{2}}\int_{0}^{1}8x(1-x)\sum_{i}c_{i}\ln\frac{L}{\Delta_{i}}\,dx = -\frac{e^{2}}{2\pi^{2}}\int_{0}^{1}x(1-x)\ln\frac{M^{2}}{\Delta}\,dx + \text{const},
\]
$\ln M^{2} = -\sum_{i\geq1}c_{i}\ln M_{i}^{2}$ at large regulator masses. The rational coefficient is $\tfrac{1}{16}\int_{0}^{1}8x(1-x)\,dx = \tfrac{1}{12}$, so the logarithm enters with $e^{2}/(12\pi^{2})$ per Dirac fermion: at large $-Q(q)$, $\Pi \approx \tfrac{e^{2}}{12\pi^{2}}\ln(-Q(q)/M^{2})$, the standard screening. For a massless chiral fermion the numerator is $2\proj$ of the same product (clause (1)); its scalar part is half the Dirac one, and its $I_{4}$ coefficient vanishes after the shift and symmetric integration (before the symmetric integration the terms linear in $\ell$ have a nonzero $I_{4}$ part, proportional to the oriented volume of $\gamma^{\mu}, \gamma^{\nu}, \ell, q$, and they integrate to zero): the pseudoscalar part of $\sum_{\alpha}\gamma^{\mu}\gamma_{\alpha}\gamma^{\nu}\gamma^{\alpha}$ and of $\gamma^{\mu}q\gamma^{\nu}q$ is zero (checked, with the control that $\langle\gamma^{\mu}\gamma_{0}\gamma^{\nu}\gamma_{1}\rangle_{4} \neq 0$ for some $\mu, \nu$). This is a statement about the numerator only: the Pauli-Villars regulators of a single Weyl fermion carry Dirac masses and break chirality, so the regulated chiral loop is not covered.

\emph{(3) Vertex.} The identities $\gamma^{\nu}a\gamma_{\nu} = -2a$, $\gamma^{\nu}ab\gamma_{\nu} = 4a\cdot b$, $\gamma^{\nu}abc\gamma_{\nu} = -2cba$ and $\gamma^{\nu}\gamma_{\nu} = 4$ hold in $\Cl(1,3)$ (checked on sample vectors; they are four-dimensional, which Pauli-Villars keeps). The one-loop vertex is $\delta\Gamma^{\mu} = -e^{2}\int\frac{d^{4}k}{(2\pi)^{4}}\gamma^{\nu}(k'+m)\gamma^{\mu}(k+m)\gamma_{\nu}\,I_{4}/(D_{k}D_{k'}D_{\gamma})$, the units $(-eI_{4})^{2}$, $I_{4}^{2}$ and $-I_{4}$ of the photon line, collected on the right, giving $-e^{2}I_{4}$. With parameters $x$ ($k$), $y$ ($k'$), $z$ (photon), $\ell = k + yq - zp$ and $\Delta = -xyQ(q) + (1-z)^{2}m^{2}$. The $Q(\ell)$ part of the numerator is $\tfrac14 Q(\ell)\sum_{\alpha}\gamma^{\nu}\gamma_{\alpha}\gamma^{\mu}\gamma^{\alpha}\gamma_{\nu} = Q(\ell)\gamma^{\mu}$ (checked), so it renormalises $F_{1}$ only. The $\ell$-free part is the numerator at $\ell = 0$. At seven points of the simplex and for two pairs of rational on-shell momenta, $p = (\tfrac54, \tfrac34, 0, 0)$, $p' = (\tfrac{13}{12}, 0, \tfrac{5}{12}, 0)$ and $p = (\tfrac53, \tfrac43, 0, 0)$, $p' = (\tfrac{17}{15}, 0, 0, \tfrac{8}{15})$, with $m = 1$, the check solves exactly for $(a, b, c)$ with $(p'+m)[\text{numerator}](p+m) = (p'+m)[a\gamma^{\mu} + b(p+p')^{\mu} + cq^{\mu}](p+m)$ for all $\mu$ and all sixteen blade components, verifies the fit exactly, and finds $b = 2mz(1-z)$ and $c = -2m(z-2)(x-y)$. Both are polynomials of degree at most two in $(y, z)$, and the seven points lie on no common conic (the six quadratic monomials have rank six on them, checked), so the identities hold on the whole simplex; that $b$ and $c$ do not depend on the external momenta is the standard reduction (Peskin and Schroeder, eq. (6.47), in their normalisation, which is $-\tfrac12$ of this numerator), checked here at two momentum pairs. The $q^{\mu}$ term is odd under $x \leftrightarrow y$, a symmetry of $\Delta$ and of the measure, so it integrates to zero, as current conservation requires. The Gordon identity $(p'+m)(p+p')^{\mu}(p+m) = (p'+m)(2m\gamma^{\mu} + \gamma^{\mu}\wedge q)(p+m)$ holds exactly, and fails with the opposite sign of the bivector (control); $\gamma^{\mu}\wedge q = -I_{5}\sigma^{\mu\nu}q_{\nu}$ in the matrix reading, since $I_{5}\sigma^{\mu\nu} = -\gamma^{\mu}\wedge\gamma^{\nu}$. With the form factors defined by $\Gamma^{\mu} = \gamma^{\mu}F_{1} + (q\wedge\gamma^{\mu})F_{2}/(2m)$, the magnetic part of $\delta\Gamma$ is $-e^{2}\cdot 2\int b(\gamma^{\mu}\wedge q)/\Delta\;I_{4}\cdot\frac{-I_{4}}{32\pi^{2}} = \frac{e^{2}}{16\pi^{2}}\int b\,(q\wedge\gamma^{\mu})/\Delta$, the two units multiplying to $-1$ (checked). So $F_{2}(0) = \frac{e^{2}}{8\pi^{2}}\,m\int_{0}^{1}dz\int_{0}^{1-z}dy\,\frac{2mz(1-z)}{(1-z)^{2}m^{2}}$. The integrand does not depend on $y$, so the inner integral gives $2z$, and $\int_{0}^{1}2z\,dz = 1$ exactly. Hence $F_{2}(0) = e^{2}/(8\pi^{2}) = \alpha/(2\pi)$ with $\alpha = e^{2}/(4\pi)$, Schwinger's value. The magnetic structure is the real bivector $\gamma^{\mu}\wedge q$: no unit enters it.

\emph{(4) Ward identity.} $(k - m)(k + m) = Q(k) - m^{2}$, so $S(k) = (k + m)/(Q(k) - m^{2})$ is the inverse of $k - m$ in $\Cl(1,3)$ off the mass shell, and $S(k+h) - S(k) = S(k+h)[(k - m) - (k + h - m)]S(k) = -S(k+h)\,h\,S(k)$, also in the other order (checked, with the controls that the opposite sign and the one-propagator form fail). Route the loop momentum $k$ through the photon line and put the regulator there; the integrands are $\Sigma(p) \propto \gamma^{\nu}S(p-k)\gamma_{\nu}$ and $\Lambda^{\mu}(p', p) \propto \gamma^{\nu}S(p'-k)\gamma^{\mu}S(p-k)\gamma_{\nu}$ with the same prefactor and photon propagator, and the identity gives $\Sigma(p+h) - \Sigma(p) = -h_{\mu}\Lambda^{\mu}(p+h, p)$ for each $k$ (checked). The regulated integrals converge after the regulator sum, so the identity survives integration: this is the Ward-Takahashi identity, and dividing by $h$ and letting $h \to 0$, $\partial\Sigma/\partial p_{\mu} = -\Lambda^{\mu}(p, p)$. With $Z_{2}^{-1} = 1 - \partial\Sigma/\partial p|_{p = m}$ and $Z_{1}^{-1} = 1 + \Lambda(p, p)|_{p = m}$ (the standard definitions), $Z_{1} = Z_{2}$.

\emph{(5) Where the unit stands.} Put $I_{4}$ for the propagator unit inside the product: the numerator becomes $4\langle\gamma^{\mu}I_{4}(k+m)\gamma^{\nu}I_{4}(k'+m)\rangle_{0}$. Since $I_{4}$ anticommutes with $\gamma^{\nu}$ and commutes with the even $k\gamma^{\nu}$, $I_{4}(k + m)\gamma^{\nu}I_{4} = -(k - m)\gamma^{\nu}$, and the numerator is $-N^{\mu\nu} + 8m^{2}g^{\mu\nu}$, $N$ the numerator of (2) (the commuting units give $-N$); checked at five $x$, three masses and two loop momenta, including $m = 0$, where the two agree. The extra $8m^{2}g$ multiplies $(\ell^{2} - \Delta)^{-2}$, whose regulated integral is $\frac{I_{4}}{16\pi^{2}}\sum c_{i}M_{i}^{2}\ln(L/\Delta_{i})$; the $\ln L$ term cancels by $\sum c_{i}M_{i}^{2} = 0$ but $-\sum c_{i}M_{i}^{2}\ln\Delta_{i}$ remains. At $q = 0$ with $M_{i}^{2} = (1, 100, 300)$ it is $-(100c_{1}\ln 100 + 300c_{2}\ln 300)$ with $c_{2} \neq 0$; a rational relation $a\ln100 + b\ln300 = 0$ means $100^{a}300^{b} = 1$, impossible for $b \neq 0$ because the prime $3$ divides $300$ and not $100$. So the remainder is not zero and the polarisation is not transverse. This matches Proposition \claimref{field-theory.second-quantisation}(4): with the unit acting as $I_{4}$ by left multiplication the two chiralities enter one complex-linearly and one conjugated, and the mass couples them, so on a massive line the unit cannot act inside the product. The unit result \claimref{synthesis.unit-i4-keeps-its-three-properties-only-in-dimension-zero-mod-four}, gives the same lesson for the regulator in $\Cl(1, D-1)$: the unit that stays covariant is the right unit, equivalently a coefficient outside the product; clause (5) is its consequence for a four-dimensional loop under Pauli-Villars.

\emph{Status of the inputs.} Clauses (1) to (5) are exact algebra in $\Cl(1,3)$ plus exact rational integrals; the covariant photon propagator on internal lines (under Premise \claimref{synthesis.krein-form-premise-indefinite-inner-product-for-gauge-fields}), the loop integrals, the parametrisation and the regulator are imported (Fact \claimref{synthesis.imported-one-loop-integrals-and-pauli-villars-regulators}; Peskin and Schroeder, An Introduction to Quantum Field Theory, 1995, sections 6.3, 7.1, 7.4 and 7.5; Pauli and Villars, Rev. Mod. Phys. 21 (1949) 434; Schwinger, Phys. Rev. 73 (1948) 416), and $\alpha/(2\pi)$ and $1/(12\pi^{2})$ are their standard values.
